\subsection{Das Elektron als geladener Massepunkt}
\label{sec:elektron-geladener-massepunkt}

Nach der Entdeckung des Elektrons stellte sich eine neue Frage:
Wie lässt sich dieses Teilchen im Rahmen der klassischen Physik beschreiben?

Die Antwort der klassischen Mechanik ist zunächst erstaunlich einfach:
Das Elektron wird als kleiner geladener Massepunkt betrachtet.
Es besitzt eine Masse, eine elektrische Ladung und bewegt sich unter dem
Einfluss äußerer Kräfte.

\[
m_e \approx 9{,}11 \cdot 10^{-31}\,\mathrm{kg},
\qquad
q_e = -e \approx -1{,}602 \cdot 10^{-19}\,\mathrm{C}.
\]

Damit unterscheidet sich das Elektron in der klassischen Beschreibung
nicht grundsätzlich von anderen Teilchen. Seine Bewegung folgt dem zweiten
Newtonschen Gesetz:

\[
\vec{F} = m_e \vec{a}.
\]

Der entscheidende Unterschied besteht jedoch darin, dass das Elektron
elektrisch geladen ist. Es reagiert deshalb auf elektrische und magnetische
Felder. Ein elektrisches Feld beschleunigt das Elektron in Richtung der
Kraftwirkung, ein magnetisches Feld lenkt seine Bahn seitlich ab.

Für ein Elektron mit der Ladung \(q_e=-e\) gilt:

\[
\vec{F} = q_e \vec{E}
\]

im elektrischen Feld und

\[
\vec{F} = q_e \, \vec{v} \times \vec{B}
\]

im magnetischen Feld. Beide Beiträge zusammen führen zur Lorentzkraft,
die im nächsten Abschnitt genauer behandelt wird.

\begin{DidacticBox}[Klassisches Bild des Elektrons]
	\label{box:klassisches-bild-des-elektrons}
	In der klassischen Physik wird das Elektron zunächst als geladener
	Massepunkt beschrieben:
	\begin{itemize}
		\item Es besitzt eine Masse \(m_e\).
		\item Es besitzt eine negative elektrische Ladung \(q_e=-e\).
		\item Seine Bewegung folgt dem zweiten Newtonschen Gesetz.
		\item Elektrische und magnetische Felder üben Kräfte auf das Elektron aus.
	\end{itemize}
	Dieses Modell ist einfach, anschaulich und technisch sehr erfolgreich.
	Es ist jedoch nicht vollständig. Viele Eigenschaften des Elektrons lassen
	sich erst mit der Quantenmechanik verstehen.
\end{DidacticBox}

Das klassische Bild ist damit ein notwendiger Zwischenschritt.
Es erklärt die Ablenkung von Elektronenstrahlen, die Bewegung in
elektrischen und magnetischen Feldern sowie viele technische Anwendungen.
Gleichzeitig bereitet es den Übergang zu einer tieferen Beschreibung vor:
Das Elektron ist nicht nur ein kleiner geladener Punkt, sondern ein
Quantenteilchen.